\chapter{Results and Discussion}
\label{ch:results}

This chapter answers the six research questions. Section~\ref{sec:champions} reports each pillar's
performance (RQ1); \S\ref{sec:research-deploy} contrasts research and
deployment metrics (RQ3); \S\ref{sec:data-algo} examines data quality versus algorithm and the
calibration ablation (RQ2, RQ4); \S\ref{sec:xai-results} presents the explainability study (RQ5);
\S\ref{sec:leakage} reports the leakage analysis (RQ6); and \S\ref{sec:context} places the numbers in
the context of the literature, followed by a discussion (\S\ref{sec:discussion}).

\textbf{A note on numbers.} All headline QWK values are reported \emph{uncalibrated} and are computed
consistently across the four pillars from saved continuous predictions; threshold calibration is
treated separately as an ablation (\S\ref{sec:data-algo}).

\section{Per-Pillar Performance (RQ1)}
\label{sec:champions}
Table~\ref{tab:champs} reports each pillar's champion cell. On \textbf{QWK the four pillars are
comparable}: classical 0.795, RNN 0.845, encoder 0.820, and LLM 0.842 fall in a narrow 0.05 band,
with the RNN highest and the classical model lowest (the classical and BiLSTM champions, which share
the same 895 test set, differ by only 0.05 QWK). In contrast, on the \textbf{deployment} metrics the
LLM is the clear leader, with the best accuracy (0.84), macro-F1 (0.84), exact match (0.67), and
within-one-point agreement (0.79).

\begin{table}[t]
\centering
\caption{Per-pillar champions (random split), standard metrics. The four QWK scores fall in a narrow
band; the LLM leads accuracy, macro-F1, Cohen $\kappa$, and the deployment metrics. Agreement columns
follow \citet{williamson2012framework}. All values uncalibrated.}
\label{tab:champs}
\small
\begin{tabular}{@{}lcccccc@{}}
\toprule
Pillar & QWK & Cohen $\kappa$ & Acc.\ & macro-F1 & Exact & Adj.\ ($\pm1$) \\
\midrule
Classical (SVR)      & 0.795 & 0.47 & 0.63 & 0.42 & 0.20 & 0.71 \\
RNN (BiLSTM)         & \textbf{0.845} & 0.62 & 0.75 & 0.53 & 0.54 & 0.77 \\
Transformer (GTE)    & 0.820 & 0.62 & 0.73 & 0.54 & 0.57 & 0.77 \\
LLM (Qwen)           & 0.842 & \textbf{0.77} & \textbf{0.84} & \textbf{0.84} & \textbf{0.67} & \textbf{0.79} \\
\bottomrule
\end{tabular}
\end{table}

The answer to RQ1 is that \emph{no single pillar dominates}. A 30-second CPU classical model is
within 0.05 QWK of a fine-tuned LLM on the research metric, while the LLM pulls clearly
ahead only on classroom deployment accuracy. These are seen-question numbers; \S\ref{sec:leakage}
shows how they change on unseen questions.

\section{Research versus Deployment (RQ3)}
\label{sec:research-deploy}
The pillars separate cleanly when research and deployment metrics are distinguished
(Table~\ref{tab:goal}). The four are comparable on QWK, so the ``winner'' there is not meaningful; but the
fine-tuned LLM wins every metric a classroom cares about: exact integer-score match (0.672), within
one point (0.788), and lowest mean absolute error ($\approx$0.98 points). On the 909 head-to-head, the
encoder and LLM tie on QWK (both 0.842) yet the LLM beats the encoder on exact match by 22.7 points.

\begin{table}[t]
\centering
\caption{Best model by goal. The QWK lead is small (the pillars are comparable); the LLM wins deployment.}
\label{tab:goal}
\small
\begin{tabular}{@{}llr@{}}
\toprule
Goal & Best pillar & Value \\
\midrule
Highest QWK (research)            & RNN \emph{(tied, n.s.)} & 0.845 \\
Highest accuracy                  & LLM & 0.839 \\
Highest exact integer match       & LLM & 0.672 \\
Highest within $\pm1$ point       & LLM & 0.788 \\
Lowest MAE (points)               & LLM & $\approx$0.98 \\
Cheapest (CPU, $\sim$30\,s)       & Classical & n/a \\
\bottomrule
\end{tabular}
\end{table}

The practical guidance is therefore to choose the tool by the goal: for the cheapest,
self-explaining option use the classical model; for classroom point-accuracy use the fine-tuned LLM.

\section{Data Quality versus Algorithm, and the Calibration Ablation (RQ2, RQ4)}
\label{sec:data-algo}
Holding the architecture fixed (the GTE dual encoder with the max-score feature), the best cell
improves $0.820\to0.842\to0.847$ in QWK across the \texttt{full}, \texttt{no10c}, and
\texttt{no10c\_no0} variants, a $+0.027$ gain from \emph{data cleaning alone}, with no change to the
algorithm. This is larger than the differences between architectures, supporting the answer to RQ2:
on small-data ASAG, data quality moves the needle at least as much as swapping models. A separate
ablation confirms robustness to the later cleaning refinement: removing the residual zero-width and
bullet noise changes the classical QWK by at most 0.004 on any variant (a negligible shift).

\paragraph{Calibration is a fragile, model-dependent lever.} Post-hoc threshold calibration helps the
\emph{classical} model (validation QWK rises $0.758\to0.829$ and test QWK $0.795\to0.847$), but
applying the \emph{same} validation-selected rule to the BiLSTM \emph{lowers} its test QWK
($0.845\to0.818$), even though its validation QWK rose. Because the four cut-points are fit on only
$\sim$134 validation rows, the gain does not transfer reliably. This is why calibration is reported as
an ablation and the headline numbers (Table~\ref{tab:champs}) are uncalibrated.

\paragraph{The max-score feature.} Supplying each question's maximum score as a normalised input
feature gives consistent gains on the deployment metrics (for example, $+4.5$ points of exact match
for the neural heads on the 895 set). It is the cheapest reliable improvement in the grid, because it
tells every model the per-question scoring scale that ``0.5'' alone cannot convey.

\section{Explainability (RQ5)}
\label{sec:xai-results}
The explainability contribution is \emph{Leave-One-Out (LOO) word attribution} (text highlighting):
the system highlights which words in the answer drove the grade using occlusion as the single
model-agnostic engine across all four families.
Table~\ref{tab:xai} reports faithfulness on 135 test answers (seed 42), and Figure~\ref{fig:faith}
shows the AOPC comprehensiveness.

\begin{table}[t]
\centering
\caption{LOO faithfulness ($n{=}135$, seed 42). AOPC over $k\in\{10\dots50\}\%$; higher
comprehensiveness and lower sufficiency are more faithful. ``Gap'' is comprehensiveness minus the
random-removal baseline ($>0$ beats random). ``Plaus.'' is the reference-overlap plausibility proxy.}
\label{tab:xai}
\small
\begin{tabular}{@{}lcccccc@{}}
\toprule
Model $\cdot$ LOO & Comp & Suff & Gap & AOPC-C & AOPC-S & Plaus. \\
\midrule
Classical $\cdot$ occlusion & 0.125 & 0.030 & $+$0.096 & \textbf{0.150} & 0.011 & 0.75 \\
BiLSTM $\cdot$ occlusion    & 0.285 & 0.082 & $+$0.274 & \textbf{0.317} & 0.043 & 0.68 \\
\bottomrule
\end{tabular}
\end{table}

\begin{figure}[t]
\centering
\includegraphics[width=0.7\textwidth]{fig_faithfulness.png}
\caption{AOPC comprehensiveness for LOO occlusion across both CPU-local families. Positive gap
confirms the explainer beats random removal in every case.}
\label{fig:faith}
\end{figure}

\textbf{LOO occlusion word attribution is reliably faithful} for every model (AOPC comprehensiveness
$+0.150$ for the classical model and $+0.317$ for the BiLSTM): removing the words it flags moves the
score far more than removing random words, and keeping only those words preserves the score (low
sufficiency). LOO works identically on the non-differentiable classical SVR and the BiLSTM,
confirming that no gradient access or internal model inspection is required. Figure~\ref{fig:heatmap}
shows an example occlusion heatmap on a real Khmer answer.

\begin{figure}[t]
\centering
\includegraphics[width=0.95\textwidth]{heatmap_classical.png}
\caption{Example Khmer word-attribution heatmap (classical, occlusion) on an original student answer;
darker means higher attribution to the predicted score. Rendered with browser text shaping, because
standard plotting back-ends do not shape Khmer complex script.}
\label{fig:heatmap}
\end{figure}

For RQ5, LOO word attribution gives a faithful, model-agnostic explanation for every pillar.
The classical model is both competitive \emph{and} self-explaining, and there is no
accuracy-versus-explainability trade-off: the unified LOO method works equally on the cheapest and
the most expensive model.

\section{Robustness to Unseen Questions (RQ6)}
\label{sec:leakage}
The random split shares questions between train and test, because only 41 distinct questions exist.
Re-evaluating the classical champion under a question-held-out split (no shared questions), across
five seeds, the mean QWK \textbf{collapses from 0.759 to 0.354} (Table~\ref{tab:leakage},
Figure~\ref{fig:leakage}). Most of the seen-question performance was therefore question-specific
memorisation rather than transferable grading skill. With only about six questions per unseen-question
test set the estimate is high-variance ($\pm0.11$), but the drop is unambiguous.

\begin{table}[t]
\centering
\caption{Classical QWK under two split modes (mean $\pm$ s.d.\ over 5 seeds).}
\label{tab:leakage}
\small
\begin{tabular}{@{}lc@{}}
\toprule
Split & Classical QWK \\
\midrule
Random (seen questions)            & $0.759 \pm 0.041$ \\
Question-held-out (unseen)         & $0.354 \pm 0.106$ \\
\midrule
Leakage gap                        & $\approx -0.40$ \\
\bottomrule
\end{tabular}
\end{table}

\begin{figure}[t]
\centering
\includegraphics[width=0.7\textwidth]{fig_leakage.png}
\caption{Question leakage: classical QWK drops from 0.76 (seen questions) to 0.35 (unseen questions).
Reported QWK should be read as seen-question performance.}
\label{fig:leakage}
\end{figure}

This is an honest headline finding: small-question-pool ASAG benchmarks, ours, and much of the
literature, are substantially inflated by question leakage. Reported QWK should be read as
\emph{seen-question} performance, and grading genuinely new questions on this corpus remains an open
problem (the neural and LLM unseen-question runs are left to future work and are expected to drop
similarly).

\section{Context: Comparison with the Literature}
\label{sec:context}
Because direct cross-dataset comparison is unsafe, two careful comparisons are made with the Arabic
anchor~\citep{alaoui2024}. First, on the \emph{same} unweighted Cohen's $\kappa$, our models score
0.47 (classical), about 0.62 (encoder and RNN), and 0.77 (LLM), versus their 0.48 (BERT) and 0.60
(best transformer); our encoder and RNN match their transformer and our LLM is higher, while our
classical matches their BERT. This is still a different dataset and class count (five-class Khmer
versus three-class Arabic), so it is context, not a leaderboard. Second, and more fairly, the
train/test \emph{generalisation gap}: their best transformer overfits by $+18.5$ points
(95.67/77.22 accuracy), whereas our deep cells are far tighter ($+1.4$ for GTE, $+2.4$ for Qwen),
and even the simple classical model ($+15.6$, 78.6/63.0) stays under their transformer. There is no
hidden overfitting beneath our reported numbers.

\section{Discussion}
\label{sec:discussion}
Taken together, the results make the ``best model'' question genuinely three-dimensional, and each
axis behaves differently. On \emph{research} agreement the pillars are comparable; on
\emph{deployment} the fine-tuned LLM is the clear winner; on \emph{explainability} LOO occlusion is reliably faithful for every model family. This picture is
consistent with the wider literature: that data quality and cheap conditioning often rival
architecture choices on small data, and that small-question-pool evaluations can over-state
generalisation. For an accountability-sensitive setting such as classroom
grading, the implication is concrete: a cheap classical model that explains itself faithfully via
occlusion is an attractive default, the LLM is the right choice when exact teacher scores matter, and
in both cases the explanation shown to a teacher should be a \emph{measured-faithful} one.
